\documentclass[11pt]{article}
\usepackage[margin=0.85in]{geometry}
\usepackage{amsmath,amssymb,amsthm,mathtools}
\usepackage{booktabs,longtable,tabularx,array,enumitem,microtype,needspace}
\usepackage{caption}
\captionsetup{font=small,labelfont=bf,skip=4pt,justification=raggedright,singlelinecheck=false}
\usepackage{xcolor}
\usepackage{hyperref}
\pdfstringdefDisableCommands{\def\ensuremath#1{#1}\def\geq{≥}}
\hypersetup{
  colorlinks=true,
  linkcolor=blue!55!black,
  citecolor=blue!55!black,
  urlcolor=blue!55!black,
  pdfauthor={Brandon Li},
  pdftitle={No binary [[14,3,d\ensuremath{\geq}5]] quantum stabilizer code exists},
  pdfsubject={A signed-shadow proof and exhaustive additive-code classification},
  pdfcreator={LaTeX with Tectonic}
}
\setlength{\parindent}{0pt}
\setlength{\parskip}{4pt plus 1pt minus 1pt}
\setlength{\emergencystretch}{1.5em}
\setcounter{secnumdepth}{3}
\newtheorem{theorem}{Theorem}
\newtheorem{lemma}[theorem]{Lemma}
\newtheorem*{corollary*}{Corollary}
\providecommand{\tightlist}{\setlength{\itemsep}{0pt}\setlength{\parskip}{0pt}}
\providecommand{\real}[1]{#1}
\newcommand{\Sh}{\mathrm{Sh}}
\newcommand{\rhs}{\mathrm{rhs}}
\title{No binary \(\left[\!\left[14,3,d\geq 5\right]\!\right]\) quantum stabilizer code exists\\[0.7em]
\large\mdseries A signed-shadow proof and exhaustive additive-code classification}
\author{Brandon Li\\[0.3em]
\normalsize Independent researcher\\
\normalsize\href{mailto:brandon.li.gn@gmail.com}{\texttt{brandon.li.gn@gmail.com}}}
\date{}
\begin{document}
\maketitle
\begin{abstract}
We prove that there is no qubit stabilizer code with parameters
\(\left[\!\left[14,3,d\geq5\right]\!\right]\). Together with the known
\(\left[\!\left[14,3,4\right]\!\right]\) code, this shows that the
largest minimum distance of a qubit stabilizer code of length 14
encoding 3 qubits is 4, which settles an open entry in the tables of
quantum code bounds. The proof has three parts. First, a signed shadow
inequality and an integrality congruence restrict the stabilizer words
of weight at most four. Second, exact Farkas certificates for split
weight enumerators rule out stabilizer words of weight two and all
configurations of three stabilizer words of weight four, so the
stabilizer contains exactly one word of weight four. Third, this word
yields an additive \((10,2^{10},\geq5)\) code with symplectic hull
dimension four. We reconstruct all 37 equivalence classes of additive
\((10,2^{10},\geq5)\) codes by exhaustive lengthening. Of these, 25
have hull dimension different from four, and each of the other 12
violates a necessary coset-capacity bound.
\end{abstract}

\noindent\textbf{Keywords:} quantum error correction; stabilizer codes; additive codes; shadow enumerators; linear programming bounds; computer-assisted proof

\noindent\textbf{Mathematics Subject Classification (2020):} 81P73, 94B65, 94B05

\section{Introduction and main result}\label{introduction-and-main-result}

\begin{theorem}[Nonexistence]\label{thm:main}
There is no binary qubit stabilizer code with parameters \(\left[\!\left[14,3,d\geq 5\right]\!\right]\).
\end{theorem}

\begin{corollary*}
The maximum distance of a binary qubit stabilizer code with \(n=14\) and \(k=3\) is \(d_{\max}(14,3)=4\).
\end{corollary*}
\begin{proof}
Theorem~\ref{thm:main} gives the upper bound, and the known
\(\left[\!\left[14,3,4\right]\!\right]\) construction gives the matching lower bound
\cite{grassl-table}.
\end{proof}

Grassl's tables of quantum codes \cite{grassl-table} list a \(\left[\!\left[14,3,4\right]\!\right]\) code and an upper bound of five on the minimum distance of a \(\left[\!\left[14,3\right]\!\right]\) qubit stabilizer code. Ball, Centelles, and Huber \cite{ball-centelles-huber} identified \(\left[\!\left[14,3,5\right]\!\right]\) in Research Problem 1 as the smallest parameter set for which the existence of a qubit stabilizer code was then unknown. Theorem~\ref{thm:main} resolves that question in the negative. Bierbrauer, Fears, Marcugini, and Pambianco \cite{bierbrauer-et-al} used a related approach to prove that no \(\left[\!\left[13,5,4\right]\!\right]\) code exists. Our argument combines shadow enumerators \cite{rains-weight,rains-shadow} with a classification of additive codes; it does not assume that the code is \(\operatorname{GF}(4)\)-linear.

The proof has three steps. Section~4 uses shadow identities to restrict the stabilizer words of low weight. Section~5 uses exact linear infeasibility certificates to rule out the remaining low-weight configurations. Sections~6 and~7 reduce the problem to additive codes of length ten, classify those codes, and finish with a coset-counting contradiction. Section~\ref{computational-material} describes the computations, and Appendix~\ref{app:systems} specifies every linear system of Section~5 in full.

Throughout, \emph{physical} weight means Pauli weight on the original fourteen qubits. The result covers both pure and impure (degenerate) stabilizer codes, but not non-stabilizer quantum codes.

\section{Pauli space and notation}\label{pauli-space-and-notation}

Let

\[
\begin{aligned}
V_{n} &=\mathbb{F}_{2}^{2n} = \{(x\mid z): x,z \in \mathbb{F}_{2}^{n}\},
\end{aligned}
\]

encode Pauli labels modulo phase. The symplectic form is

\[
\begin{aligned}
\langle (x\mid z),(x'\mid z') \rangle &=x \cdot z' + z \cdot x' \in \mathbb{F}_{2}.
\end{aligned}
\]

The Pauli weight is

\[
\begin{aligned}
\operatorname{wt}(x\mid z) &=\#\{i : (x_{i},z_{i}) \ne (0,0)\}.
\end{aligned}
\]

A binary stabilizer code with parameters \(\left[\!\left[14,3,d\right]\!\right]\) is an isotropic subspace \(S \leq V_{14}\) of binary dimension \(14-3=11\). Its centralizer is \(S^\perp\), and its distance is

\[
\begin{aligned}
d(S) &=\min \{ \operatorname{wt}(v) : v \in S^\perp \setminus S \}.
\end{aligned}
\]

We call each element of \(S\) a stabilizer word. A stabilizer generator is an element of a chosen basis of \(S\); weight enumerators count all stabilizer words, not only basis generators.

For any nonzero binary subspace \(C\) of a Pauli space, write
\[
d(C)=\min\{\operatorname{wt}(c):c\in C\setminus\{0\}\}.
\]
We write \(d_{\max}(n,k)\) for the greatest minimum distance \(d(S)\) among binary stabilizer codes with fixed \(n\) and \(k\).

We assume for contradiction that \(d(S) \geq 5\). A code is \emph{pure to distance} \(d\) if no nonzero stabilizer word has weight below \(d\); otherwise it is degenerate (or impure). Since the distance is taken over \(S^\perp\setminus S\), stabilizer words of low weight are allowed. Under our assumption, every vector of weight at most four that commutes with \(S\) lies in \(S\).

For a binary subspace \(C\) of a Pauli space, its symplectic hull is \(\operatorname{Hull}(C)=C \cap C^\perp\). When a block of coordinates is specified, \(H_{a,b}\) below counts stabilizer words of split weight \((a,b)\); it is an enumerator, not a subspace. The hull in Section~6 is denoted \(J\) to keep the two notions distinct.

\section{The signed affine shadow and split weight enumerators}\label{the-signed-affine-shadow-including-split-weights}

For one qubit define

\[
\begin{aligned}
q(x,z) &=x + z + xz \in \mathbb{F}_{2}.
\end{aligned}
\]

This is one exactly when \((x,z)\) is nonzero, so on \(V_{n}\) it satisfies

\[
\begin{aligned}
q(v) &\equiv\operatorname{wt}(v) \pmod 2, \\
q(v+w) &=q(v) + q(w) + \langle v,w \rangle.
\end{aligned}
\]

Because \(S\) is isotropic, \(q|_{S}\) is linear. Nondegeneracy of the symplectic form gives a vector \(v_{0}\) such that

\[
\begin{aligned}
\langle v_{0},s \rangle &=q(s)\quad\forall s \in S.
\end{aligned}
\]

Character orthogonality gives the indicator identity

\[
\begin{aligned}
1_{v_{0}+S^\perp}(v)
& = 2^{-11} \sum_{s \in S} (-1)^{q(s)+\langle v,s \rangle}.
\end{aligned}
\]

Thus the signed transform counts the vectors of each weight in the affine coset \(v_{0}+S^\perp\), the shadow of \(S\), so every entry of the split shadow enumerator below is a nonnegative integer.

Fix a partition of the fourteen qubits into blocks of lengths \(m\) and \(f\), with \(m+f=14\), and let \(H_{i,j}\) count stabilizer words with block weights \((i,j)\). Define the quaternary Krawtchouk coefficients

\[
\begin{aligned}
K_{a}^{m}(i) &=[u^{a}] (1+3u)^{m-i} (1-u)^{i}.
\end{aligned}
\]

If \(C_{a,b}\) counts \(S^\perp\) and \(\Sh_{a,b}\) counts the affine shadow, then

\[
\begin{aligned}
2048 C_{a,b}
& = \sum_{i,j} H_{i,j} K_{a}^{m}(i) K_{b}^{f}(j), \\[0.4ex]
2048 \Sh_{a,b}
& = \sum_{i,j} (-1)^{i+j} H_{i,j} K_{a}^{m}(i) K_{b}^{f}(j).
\end{aligned}
\]

We write \(T_{a,b}=2048C_{a,b}\) and \(\mathtt{SH}_{a,b}=2048\Sh_{a,b}\) for the right-hand sides; these are integer linear forms in the \(H_{i,j}\).

Both \(C_{a,b}\) and \(\Sh_{a,b}\) are nonnegative integers, and \(C_{a,b} \geq H_{a,b}\) because \(S \leq S^\perp\). If \(1 \leq a+b \leq 4\), then

\[
\begin{aligned}
C_{a,b} &=H_{a,b},
\end{aligned}
\]

because every vector of \(S^\perp\) of weight at most four lies in \(S\) when \(d(S) \geq 5\). This holds whether or not \(S\) is pure.

\section{Global low-weight restrictions and parity}\label{global-low-weight-restrictions-and-parity}

Write \(A_{j}\) for the number of stabilizer words of total physical weight \(j\), and put

\[
\begin{aligned}
T_{j} &=\sum_{w} A_{w} K_{j}^{14}(w), \\
U_{j} &=\sum_{w} (-1)^{w} A_{w} K_{j}^{14}(w).
\end{aligned}
\]

Write \(\Sh_j=\sum_{a+b=j}\Sh_{a,b}\) for the total-weight count in the affine shadow. The quantities \(T_{j}/2048\) and \(U_{j}/2048\) are the actual total-weight counts in \(S^\perp\) and in the affine shadow. For the fixed partition of Section~3, coefficient convolution gives

\[
\begin{aligned}
\sum_{a+b=j}K_a^m(i)K_b^f(\ell)
&=[t^j](1+3t)^{m-i}(1-t)^i(1+3t)^{f-\ell}(1-t)^\ell\\
&=K_j^{14}(i+\ell).
\end{aligned}
\]

Summing the split-shadow transform of Section~3 over \(a+b=j\) therefore gives
\[
U_j=\sum_{a+b=j}\mathtt{SH}_{a,b}=2048\Sh_j.
\]

Hence \(U_{j} \geq 0\) for every \(j\). Also \(A_{0}=1\), \(\sum_{j} A_{j}=2048\), and \(T_{j}=2048 A_{j}\) for \(j=1,2,3,4\).

\begin{lemma}[Low-weight restriction]\label{lem:low-weight}
If \(d(S)\geq5\), then \(A_1=A_3=0\), \(A_2\leq2\), and \(A_4\leq3\).
\end{lemma}
\begin{proof}
The proof uses only the following identity between the \(A_j\), \(T_j\), and \(U_j\); its coefficients were found by computer, and it can be checked by expanding the Krawtchouk transforms in \(T_{j}\) and \(U_{j}\):

\[
\begin{aligned}
989184 A_{4} &=-296 U_{2} -88 U_{4} -9584640 A_{0} +6309 \sum_{j} A_{j} \\
& -11685888 A_{1} -1345536 A_{2} -4681728 A_{3} \\
& +1914 (T_{1}-2048 A_{1}) +1318 (T_{2}-2048 A_{2}) \\
& +279 (T_{3}-2048 A_{3}) +161 (T_{4}-2048 A_{4}) \\
& -1081344 A_{5} -254976 A_{6} -15360 A_{7} -44032 A_{14}.
\end{aligned}
\]

Now \(A_0=1\), \(\sum_j A_j=2048\), and \(T_j=2048A_j\) for \(1\leq j\leq4\). Since \(U_2,U_4,A_5,A_6,A_7,A_{14}\geq0\), the identity implies

\[
\begin{aligned}
11685888 A_{1} +1345536 A_{2} +4681728 A_{3} +989184 A_{4} &\leq 3336192.
\end{aligned}
\]

Since the \(A_j\) are nonnegative integers, these coefficients imply

\[
\begin{aligned}
A_{1} &=A_{3} = 0,\qquad A_{2} \leq 2,\qquad A_{4} \leq 3.
\end{aligned}
\]
\end{proof}

\noindent\emph{Remark.} Section~5.3 excludes the remaining cases \(A_{2}\in\{1,2\}\).

\subsection{Integral affine-shadow parity}\label{integral-affine-shadow-parity}

Lemma~\ref{lem:low-weight} does not rule out \(A_{4}=0\), but an integrality congruence does. The coefficients below are chosen so that, after expanding the MacWilliams and shadow equations, only the terms in \(A_{3}\) and \(A_{4}\) survive modulo 4096. Let

\[
\begin{aligned}
\Delta_{0} &=\sum_{w} A_{w}, \\
\Delta_{j} &=T_{j} -2048 A_{j},\qquad j=1,\dots,4, \\
\Gamma_{j} &=U_{j} -2048 \Sh_{j},\qquad j=0,\dots,3.
\end{aligned}
\]

Take the integer combination

\[
\begin{aligned}
L &=\Delta_{0} -2\Delta_{1} -6\Delta_{2} -17\Delta_{3} -95\Delta_{4} \\
& -20\Gamma_{0} -24\Gamma_{1} +1108\Gamma_{2} +246\Gamma_{3}.
\end{aligned}
\]

Since \(\Delta_0=2048\) and \(\Delta_j=\Gamma_j=0\) for the other indices, \(L=2048\). Writing \(c_j\) for the coefficient of \(A_j\), direct expansion gives

\[
\begin{aligned}
(c_0,\dots,c_{14})={}&(-4550656,-7348224,-1052672,-2623488,321536,\\
&\qquad -634880,221184,-28672,24576,24576,\\
&\qquad -45056,0,40960,-16384,-77824).
\end{aligned}
\]

Modulo 4096 every coefficient vanishes except those of \(A_{3}\) and \(A_{4}\), each congruent to 2048. The four shadow-variable coefficients are \((40960, 49152, -2269184, -503808)\), all divisible by 4096. Since \(L=2048\), reduction modulo 4096 gives
\[
A_{3}+A_{4}\equiv1\pmod 2.
\]

Together with \(A_{3}=0\) and \(0 \leq A_{4} \leq 3\), this gives

\[
\begin{aligned}
A_{4} \in \{1,3\}.
\end{aligned}
\]

\section{Excluding low-weight stabilizer configurations}\label{excluding-low-weight-stabilizer-configurations}

For each subgroup \(R\) of low-weight stabilizer words that is still possible, we write down a linear system that every code containing \(R\) must satisfy, and give a Farkas certificate showing that the system has no solution.

No step uses a later one. Section~5.3 proves \(A_2=0\) using only Lemma~\ref{lem:low-weight}. Section~5.4 rules out certain pairs of weight-four words, and Section~5.5 uses this to list the possible configurations when \(A_4=3\). Sections~5.6 and~5.7 rule out each of these configurations, using \(A_2=0\). Lemma~\ref{lem:unique-check} combines the results.

\subsection{A necessary coset-counting relaxation}\label{a-necessary-coset-counting-relaxation}

Let \(R \leq S\) be an isotropic subgroup supported on \(m\) of the original qubits, with the remaining \(f=14-m\) sites forming a suffix. Every element of \(S\) commutes with \(R\), so its prefix lies in \(R^\perp\). Also, \(S\) is a union of complete \(R\)-cosets.

For each complete coset \(Q\) of \(R\) in \(R^\perp\), let

\[
\begin{aligned}
p(Q)&=(p_{0},\dots,p_{m}), \\
  p_{a} &=\#\{q \in Q : \operatorname{wt}_{\mathrm{pref}}(q)=a\},\qquad 0\leq a\leq m.
\end{aligned}
\]

Here \(\operatorname{wt}_{\mathrm{pref}}\) is Pauli weight on the original prefix sites.
Let \(y_{p,b}\) be the number of complete \(R\)-cosets selected by \(S\), aggregated over all suffix Pauli labels of weight \(b\), and grouped only by the prefix pattern \(p\). Then \(y_{p,b} \geq 0\) and

\[
\begin{aligned}
H_{a,b} &=\sum_{p} p_{a} y_{p,b}.
\end{aligned}
\]

The model requires only \(R\leq S\); it does not assume that \(R\) is all of \(S\cap V_{\mathrm{pref}}\). If \(S\) contains further words supported on the prefix, they are still counted correctly, because \(S\) is a union of complete \(R\)-cosets. Every code containing \(R\) therefore gives a solution. The model forgets which cosets and suffix labels occur and how different suffixes fit together linearly, so it may also have solutions that come from no code. Infeasibility therefore excludes the code, while feasibility would prove nothing. Appendix~\ref{app:systems} defines every row and column of each system exactly and lists the coset patterns.

For every such model we use the following row families:

\begin{itemize}
\tightlist
\item
  \(H_{0,0}=1\) and \(\sum_{a,b}H_{a,b}=2048\);
\item
  physical coset equalities or inequalities forced by the selected subgroup;
\item
  the low-weight distance equations \(T_{a,b}-2048H_{a,b}=0\) for \(1 \leq a+b \leq 4\);
\item
  centralizer inclusion \(2048H_{a,b}-T_{a,b} \leq 0\) in every split cell;
\item
  centralizer nonnegativity \(-T_{a,b} \leq 0\) in every split cell;
\item
  signed affine-shadow nonnegativity \(-\mathtt{SH}_{a,b} \leq 0\) in every split cell; and
\item
  \(y \geq 0\).
\end{itemize}

Sections~5.6 and~5.7 also use the global weight equations stated there. The coset-capacity bound of Section~6.2 is not used in Section~5.

\subsection{Exact Farkas sign convention}\label{exact-farkas-sign-convention}

After choosing the variables for a particular relaxation and scaling by 2048, write the necessary system as

\[
\begin{aligned}
M z &=\rhs,\qquad G z \leq 0,\qquad z \geq 0.
\end{aligned}
\]

A Farkas certificate consists of nonnegative \(\lambda\), free-sign \(\nu\), and a coefficient vector \(c\) satisfying

\[
\begin{aligned}
\lambda &\geq 0, \\
c &=\lambda G + \nu M, \qquad c\geq 0, \\
\nu \cdot \rhs &< 0.
\end{aligned}
\]
The inequality \(c\geq0\) is entrywise.

For a putative feasible \(z\),

\[
\begin{aligned}
0 \leq c \cdot z &=\lambda \cdot (G z) + \nu \cdot (M z) \\
 &\leq \nu \cdot \rhs < 0,
\end{aligned}
\]

which is impossible. Every certificate is checked in exact integer arithmetic.

\subsection{Excluding weight-two stabilizers}\label{excluding-weight-two-stabilizers}

Suppose \(p=Z_{0} Z_{1}\) belongs to \(S\). This loses no generality: single-qubit Clifford operations permute the three nonidentity Pauli matrices transitively and preserve weight, so a local Clifford operation on sites 0 and 1 maps any weight-two stabilizer word to \(Z_0Z_1\) without changing the code parameters. The eight two-site Pauli labels commuting with \(p\) split into four complete cosets of \(\langle p \rangle\) with two-site pattern

\[
\begin{aligned}
(1,0,1)&:1, \qquad (0,2,0):1, \qquad (0,0,2):2.
\end{aligned}
\]
The integers after the colons are the coset multiplicities.

For every fixed suffix Pauli label, complete-coset closure therefore gives

\[
\begin{aligned}
H_{2,b} &\geq H_{0,b},\qquad b=0,\dots,12.
\end{aligned}
\]

The 39 variables are the split counts \(H_{a,b}\) of stabilizer words with prefix weight \(a\in\{0,1,2\}\) on sites 0 and 1 and suffix weight \(b\in\{0,\ldots,12\}\) on the remaining sites.

\Needspace{10\baselineskip}
Every code containing \(Z_0Z_1\) satisfies the following constraints:

\begin{itemize}
\item \(H_{0,0}=1\), \(\sum_{a,b}H_{a,b}=2048\), and \(H_{a,b}\geq0\);
\item the 11 equalities \(T_{a,b}=2048H_{a,b}\) for \(1\leq a+b\leq4\);
\item the 13 coset inequalities \(H_{2,b}\geq H_{0,b}\); and
\item in each of the 39 split cells, centralizer inclusion \(T_{a,b}\geq2048H_{a,b}\), centralizer nonnegativity \(T_{a,b}\geq0\), and affine-shadow nonnegativity \(\mathtt{SH}_{a,b}\geq0\).
\end{itemize}

Together with the two normalization equalities, these give 13 equality rows. The inequalities give 130 rows, in addition to the 39 nonnegativity conditions on the variables. This relaxation omits the further coset-pattern relations in Section~5.1.

The Farkas certificate has 13 nonnegative inequality multipliers and 10 nonzero equality multipliers. The residual vector \(c\) has 17 positive and 22 zero entries and satisfies the 39 column identities \(c=\lambda G+\nu M\). Its normalization is

\[
\begin{aligned}
\nu \cdot \rhs &=-2{,}571{,}108{,}352.
\end{aligned}
\]

Hence \(A_{2}=0\).

\subsection{Excluding overlapping low-weight stabilizer words}\label{excluding-overlapping-low-weight-checks}

Section~5.5 uses the following exclusions. In each case the cosets of the subgroup in its normalizer are enumerated directly; Appendix~\ref{app:systems} lists the resulting patterns.

\subsubsection{A weight-three subgroup}\label{a-weight-three-subgroup}

For \(R=\langle \mathtt{ZZZ} \rangle\) on three original sites, the 16 normalizer cosets have prefix patterns

\[
\begin{aligned}
(0,0,1,1)&:12, \\
(0,1,1,0)&:3, \\
(1,0,0,1)&:1.
\end{aligned}
\]
As above, the integer after each colon is the number of cosets with that prefix-weight pattern.

The Farkas certificate satisfies the sign conditions of Section~5.2 on all 36 columns and has \(\nu \cdot \rhs=-22{,}528\). This gives a second proof that \(A_3=0\) that does not use Lemma~\ref{lem:low-weight}.

\subsubsection{The four-site Bell subgroup}\label{the-four-site-bell-subgroup}

For \(R=\langle \mathtt{XXXX},\mathtt{ZZZZ} \rangle\) on four original sites, the 16 complete cosets have patterns

\[
\begin{aligned}
(1,0,0,0,3)&:1, \\
(0,0,2,0,2)&:9, \\
(0,0,0,4,0)&:6.
\end{aligned}
\]
The integers after the colons are the coset multiplicities.

For every suffix weight \(b\),

\[
\begin{aligned}
H_{1,b}&=0, \\
H_{4,b}&=H_{2,b}+3H_{0,b}.
\end{aligned}
\]

\Needspace{4\baselineskip}
The Farkas certificate satisfies the sign conditions of Section~5.2 on all 55 columns and has

\[
\begin{aligned}
\nu \cdot \rhs &=-320{,}017{,}816{,}092{,}672.
\end{aligned}
\]

This excludes a pair of weight-four Bell stabilizer words with the same support.

\subsubsection{The five-site overlap subgroup}\label{the-five-site-overlap-subgroup}

If two commuting weight-four words overlap in three sites and differ on two of those sites, choose a local Clifford operation independently at each site. At each of the two sites where the letters differ, map the ordered pair of letters to \((X,Z)\). At the third common site, and at each site used by only one word, map the letter to \(Z\). After a site permutation this gives

\[
\begin{aligned}
p &=X_{0} X_{1} Z_{2} Z_{3} I_{4}, \\
q &=Z_{0} Z_{1} Z_{2} I_{3} Z_{4}.
\end{aligned}
\]

All three nonzero words in \(\langle p,q \rangle\) have weight four on five sites. The normalizer has 64 four-element cosets and nine prefix-weight patterns. The Farkas certificate satisfies the sign conditions of Section~5.2 on all 90 columns and has

\[
\begin{aligned}
\nu \cdot \rhs &=-52{,}501{,}014{,}528.
\end{aligned}
\]

\subsubsection{The two six-site overlap subgroups}\label{the-two-six-site-overlap-subgroups}

For the same-letter two-site overlap, use

\[
\begin{aligned}
R &=\langle \mathtt{ZZZZII}, \mathtt{ZZIIZZ} \rangle.
\end{aligned}
\]

The six-site normalizer has 256 four-element cosets and 13 prefix patterns. The corresponding 117-variable system has a Farkas certificate satisfying the sign conditions of Section~5.2 on all 117 columns, with \(\nu \cdot \rhs=-8{,}820{,}736\).

For the crossed two-site overlap, use

\[
\begin{aligned}
R &=\langle \mathtt{ZZZZII}, \mathtt{XXIIXX} \rangle.
\end{aligned}
\]

The nonzero subgroup words have weights \(4,4,6\). Here the normalizer has 14 prefix patterns; the resulting 126-variable system has a Farkas certificate satisfying the sign conditions of Section~5.2 on all 126 columns, with \(\nu \cdot \rhs=-22{,}212{,}608\).

\subsection{Support geometry of weight-four stabilizer words}\label{support-geometry-of-weight-four-checks}

Let \(p,q\) be distinct commuting weight-four stabilizer words. Let \(t\) be the number of common support sites and \(h\) the number of common sites on which their nonidentity Pauli letters differ. Commutation makes \(h\) even, and

\[
\begin{aligned}
\operatorname{wt}(p+q) &=8 - 2t + h.
\end{aligned}
\]

The preceding subgroup exclusions remove every case with \(t \geq 2\), as shown in Table~\ref{tab:obstruction}:

{\small
\setlength{\tabcolsep}{8pt}
\renewcommand{\arraystretch}{1.08}
\begin{center}
\captionof{table}{Subgroup exclusions covering every pair of distinct commuting weight-four stabilizer words that overlap in \(t\geq2\) sites, where \(h\) counts the common sites carrying different Pauli letters.}
\label{tab:obstruction}
\begin{tabular}{@{}rrl@{}}
\toprule
\(t\) & \(h\) & obstruction \\
\midrule
4 & 2 & weight-two stabilizer \\
4 & 4 & four-site Bell subgroup \\
3 & 0 & weight-two stabilizer \\
3 & 2 & five-site overlap subgroup \\
2 & 0 & six-site same-letter triangle subgroup \\
2 & 2 & six-site crossed-overlap subgroup \\
\bottomrule
\end{tabular}
\end{center}
}

The omitted case \(t=4,h=0\) would give the same Pauli word twice, not two distinct stabilizer words. Thus two distinct weight-four stabilizer words meet in at most one qubit. If they meet in one, commutation forces the same local Pauli letter there. Hence every qubit carries at most one nonidentity letter across all weight-four stabilizer words, and a single local Clifford operation at each qubit sends that letter to \(Z\). Local Clifford transformations therefore map all weight-four stabilizer words to Z-only words simultaneously. This does not imply that the full stabilizer is CSS.

\begin{minipage}{\linewidth}
If \(A_{4}=3\), the three weight-four stabilizer words are linearly independent: the sum of any two has weight six or eight, not four. Their support sets are pairwise disjoint or meet in one site. Up to site permutation and reordering, there are exactly five support patterns:

{\small
\setlength{\tabcolsep}{4pt}
\captionof{table}{The five support patterns for three weight-four stabilizer words when \(A_{4}=3\), up to site permutation and reordering. Only the first row has pairwise disjoint supports.}
\label{tab:support-geometry}
\begin{tabularx}{\linewidth}{@{}>{\raggedright\arraybackslash}p{2.7cm}Xr>{\centering\arraybackslash}p{2.1cm}r@{}}
\toprule
type & word supports & union & pair-product weights & triple-product weight \\
\midrule
0 intersections & \(\{0,1,2,3\};\{4,5,6,7\};\{8,9,10,11\}\) & 12 & \(8,8,8\) & 12 \\
1 intersection & \(\{0,1,2,3\};\{3,4,5,6\};\{7,8,9,10\}\) & 11 & \(6,8,8\) & 10 \\
2 intersections & \(\{0,1,2,3\};\{3,4,5,6\};\{6,7,8,9\}\) & 10 & \(6,8,6\) & 8 \\
3 distinct intersections & \(\{0,1,2,3\};\{3,4,5,6\};\{2,4,7,8\}\) & 9 & \(6,6,6\) & 6 \\
one common site & \(\{0,1,2,3\};\{3,4,5,6\};\{3,7,8,9\}\) & 10 & \(6,6,6\) & 10 \\
\bottomrule
\end{tabularx}
}
\end{minipage}

To see that the list is complete, record which of the three pairs of supports intersect. There can be zero, one, two, or three intersecting pairs. When exactly two pairs intersect, their intersection sites must be distinct: if both were the same, that site would lie in all three supports and the third pair would intersect too. When all three pairs intersect, either their three intersection sites differ or all three supports share one site. These are precisely the five cases above.

\subsection{The four overlapping rank-three geometries}\label{the-four-overlapping-rank-three-geometries}

For each of the four non-disjoint rows of Table~\ref{tab:support-geometry}, let \(R\) be the rank-three subgroup generated by the three displayed Z-only words. Its cosets in the normalizer are grouped by prefix-weight pattern as in Section~5.1. In each case the system also includes the global equations

\[
\begin{aligned}
\sum_{a+b=j} H_{a,b} &=0,\qquad j=1,2,3, \\
\sum_{a+b=4} H_{a,b} &=3.
\end{aligned}
\]

The first three hold by Lemma~\ref{lem:low-weight} and Section~5.3; the last is the case assumption \(A_{4}=3\). Table~\ref{tab:rank-three} summarizes the four systems and their certificates.

{\small
\setlength{\tabcolsep}{3pt}
\captionof{table}{Farkas certificates excluding the four non-disjoint geometries of Table~\ref{tab:support-geometry}. The fifth column counts the inequality rows with a positive multiplier; Table~\ref{tab:app-systems} also counts the equality rows used.}
\label{tab:rank-three}
\begin{tabularx}{\linewidth}{@{}X>{\centering\arraybackslash}p{1.0cm}>{\centering\arraybackslash}p{1.4cm}>{\centering\arraybackslash}p{1.6cm}>{\centering\arraybackslash}p{1.9cm}>{\centering\arraybackslash}p{2.1cm}@{}}
\toprule
geometry & sites & patterns & variables & rows with \(\lambda>0\) & \(\nu\cdot\rhs\) \\
\midrule
three distinct intersections & 9+5 & 40 & 240 & 12 & \(-86{,}470{,}656\) \\
one common site & 10+4 & 27 & 135 & 13 & \(-9{,}031{,}680\) \\
two intersections & 10+4 & 60 & 300 & 14 & \(-20{,}484{,}096\) \\
one intersection & 11+3 & 74 & 296 & 14 & \(-61{,}440\) \\
\bottomrule
\end{tabularx}
}

In each case the certificate satisfies the sign conditions of Section~5.2 in every column, with the value of \(\nu\cdot\rhs\) shown, so the system is infeasible.

\Needspace{8\baselineskip}
\subsection{\texorpdfstring{The disjoint \(4+4+4+2\) geometry}{The disjoint 4+4+4+2 geometry}}\label{the-disjoint-4442-geometry}

The only configuration with \(A_{4}=3\) left by Section~5.6 is

\[
\begin{aligned}
R &=\langle \mathtt{ZZZZIIIIIIIIII}, \mathtt{IIIIZZZZIIIIII}, \mathtt{IIIIIIIIZZZZII} \rangle.
\end{aligned}
\]

Each four-site block has 64 cosets of \(\langle \mathtt{ZZZZ} \rangle\) in its even-X-parity normalizer. The six weight patterns of Table~\ref{tab:block-patterns} have multiplicities totaling \(1+4+3+24+24+8=64\):

{
\begin{longtable}[]{@{}lr@{}}
\caption{Prefix-weight patterns of the 64 cosets of \(\langle \mathtt{ZZZZ} \rangle\) in one four-site block.}
\label{tab:block-patterns}\\
\toprule\noalign{}
pattern \((p_{0},p_{1},p_{2},p_{3},p_{4})\) & multiplicity \\
\midrule\noalign{}
\endhead
\bottomrule\noalign{}
\endlastfoot
\((1,0,0,0,1)\) & 1 \\
\((0,1,0,1,0)\) & 4 \\
\((0,0,2,0,0)\) & 3 \\
\((0,0,1,0,1)\) & 24 \\
\((0,0,0,2,0)\) & 24 \\
\((0,0,0,0,2)\) & 8 \\
\end{longtable}
}

Let \(y_{p,q,r,b}\) count the complete \(R\)-cosets in \(S\) whose three blocks have patterns \(p,q,r\) and whose last two sites have weight \(b\). There are \(6\cdot6\cdot6\cdot3=648\) nonnegative variables. As in Section~5.1, this is a relaxation.

The four-block split enumerator has \(5^3\cdot3=375\) cells, indexed by the three four-site block weights and the weight on the two untouched sites, and the system has 70 equalities. A Farkas certificate using the signed-shadow nonnegativity constraints gives

\[
\begin{aligned}
\lambda &\geq 0, \\
c &=\lambda G + \nu M \geq 0, \\
\nu \cdot \rhs &=-125{,}829{,}120.
\end{aligned}
\]

The equality rows include the global weight equations \(A_{1}=A_{2}=A_{3}=0\) and \(A_{4}=3\), as in Section~5.6. All 648 residual coefficients are nonnegative (403 positive and 245 zero), so this configuration is also impossible.

\begin{lemma}[Unique weight-four word]\label{lem:unique-check}
Every binary \(\left[\!\left[14,3,d\geq5\right]\!\right]\) stabilizer code has exactly one stabilizer word of weight four.
\end{lemma}
\begin{proof}
The parity restriction gives \(A_4\in\{1,3\}\), and Sections~5.3--5.7 exclude \(A_4=3\). Hence \(A_4=1\).
\end{proof}

\section{Reduction to a ten-site additive code}\label{reduction-to-a-ten-site-additive-code}

When the weight-four word \(p\) is the only stabilizer word of weight at most four, the rest of the stabilizer is determined by its restriction to the other ten qubits together with a map to the four sites of \(p\). The lemma below turns this into a statement about a single additive code of length ten.

\Needspace{11\baselineskip}
\begin{lemma}[Ten-site reduction]\label{lem:ten-site}
Suppose \(S\leq V_{14}\) is isotropic of dimension 11, has \(d(S)\geq5\), and its only nonzero stabilizer word of physical weight at most four is a word \(p\) of weight four. Equivalently, \(A_1=A_2=A_3=0\) and \(A_4=1\). Then there exists a binary subspace \(C\leq V_{10}\) of dimension 10, minimum Pauli weight at least five, and \(\dim(C\cap C^\perp)=4\).
\end{lemma}
\begin{proof}
Let \(p\) be the unique weight-four stabilizer word. After local Clifford conjugation and a site permutation, assume \(p=Z_{0}Z_{1}Z_{2}Z_{3}\) on the first four original qubits. Write \(V_{4}\) for their Pauli space and \(W_{10}\) for the remaining ten-qubit Pauli space, so \(V_{14}=V_{4} \oplus W_{10}\).

Define the six-dimensional quotient

\[
\begin{aligned}
E &=(p^\perp \cap V_{4}) / \langle p \rangle.
\end{aligned}
\]

The quotient is nondegenerate symplectic: \(p^\perp \cap V_{4}\) has binary dimension seven, and quotienting by its one-dimensional radical \(\langle p \rangle\) leaves dimension six.

By the hypothesis, \(S \cap V_{4}=\langle p \rangle\), so projection of \(S/\langle p \rangle\) to \(W_{10}\) is injective. Let its image be a binary ten-space \(P \leq W_{10}\). Then there is a linear map \(\tau:P\to E\) such that

\[
\begin{aligned}
S/\langle p \rangle &=\{ \tau(x)+x : x \in P \}, \\
\langle x,y \rangle_{W} &=\langle \tau(x),\tau(y) \rangle_{E}.
\end{aligned}
\]

\subsection{Distance forces surjectivity}\label{distance-forces-surjectivity}

Suppose \(z \in E\) is orthogonal to \(\operatorname{im}(\tau)\). Choose a four-site representative \(v\) of \(z\) in \(p^\perp\). It commutes with \(p\) and with every element of the graph, hence \(v \in S^\perp\). Its weight on the original qubits is at most four. Distance at least five forces \(v \in S\). But \(S \cap V_{4}=\langle p \rangle\), so \(z=0\). Thus \((\operatorname{im}(\tau))^\perp=0\). Nondegeneracy of \(E\) gives \(\operatorname{im}(\tau)=E\), and therefore

\[
\begin{aligned}
\operatorname{rank}(\tau)=6,\qquad \dim \operatorname{ker}(\tau)=10-6=4.
\end{aligned}
\]

Surjectivity also identifies the radical of the pairing on \(P\):

\[
\begin{aligned}
\operatorname{rad}(P)=P \cap P^\perp_{W}=\operatorname{ker}(\tau),\qquad \dim \operatorname{rad}(P)=4.
\end{aligned}
\]

Put

\[
\begin{aligned}
C &=P^\perp_{W}.
\end{aligned}
\]

Then \(\dim C=10\), so \(\lvert C\rvert=1024\). If a nonzero \(c \in C\) had physical weight at most four, the suffix-only vector \((0,c)\) would be in \(S^\perp\), and distance would force it into \(S\). It cannot equal \(p\), since its prefix is zero; it would therefore be another nonzero stabilizer word of weight at most four, contrary to the hypothesis. Thus \(d(C) \geq 5\). Because \(C^\perp=P\),

\[
\begin{aligned}
\dim(C \cap C^\perp_{W})=\dim(P \cap C)=4.
\end{aligned}
\]
\end{proof}

\subsection{A necessary coset-capacity bound}\label{a-necessary-coset-capacity-bound}

This subsection proves a counting bound on \(C\). Each nonzero coset \(\xi\) defined below is matched with a different nonzero class \(e(\xi)\) of four-site prefixes. A centralizer element built from the two has weight \(d(\xi)+\operatorname{cost}(e(\xi))\), which must be at least five, so low-weight cosets need high-cost prefix classes. There are only 8, 32, and 59 prefix classes of cost at least 4, 3, and 2, and this limits the number of low-weight cosets.


Let \(J=C \cap C^\perp_{W}\), so \(\dim J=4\), and let \(P=C^\perp_{W}\). The quotient \(P/J\) is a nondegenerate six-dimensional symplectic space. Also \(J^\perp_{W}/C\) has 64 suffix cosets. For \(\xi=w+C\in J^\perp_{W}/C\), put
\[
d(\xi)=\min_{c\in C}\operatorname{wt}(w+c).
\]
\begin{minipage}{\linewidth}
The pairing
\[
(J^\perp_{W}/C)\times(P/J)\longrightarrow\mathbb{F}_{2},
\qquad
(w+C,x+J)\longmapsto\langle w,x\rangle_{W}
\]
is well defined and nondegenerate: its left kernel is \(P^\perp_{W}=C\), and both quotient spaces have dimension six.
\end{minipage}

For each \(\xi=w+C\), there is a unique \(e(\xi)\in E\) such that
\[
\langle e(\xi),\tau(x)\rangle_{E}=\langle w,x\rangle_{W}
\qquad(x\in P).
\]
The symplectic form on \(P/J\) identifies \((P/J)^*\) with \(P/J\), and \(\tau\) induces a symplectic isomorphism \(P/J\to E\). The perfect pairing above therefore identifies \(J^\perp_{W}/C\) with \(E\), so \(\xi\mapsto e(\xi)\) is a bijection and sends zero to zero.

For \(e\in E\), define its physical prefix cost by
\[
\operatorname{cost}(e)=
\min\{\operatorname{wt}(v):v\in p^\perp\cap V_{4},\ v+\langle p\rangle=e\}.
\]
Its distribution is given in Table~\ref{tab:cost}.

{\small
\setlength{\tabcolsep}{8pt}
\renewcommand{\arraystretch}{1.08}
\begin{center}
\captionof{table}{Distribution of the physical prefix cost over the 64 classes of \(E\).}
\label{tab:cost}
\begin{tabular}{@{}lccccc@{}}
\toprule
cost & 0 & 1 & 2 & 3 & 4 \\
\midrule
number of classes & 1 & 4 & 27 & 24 & 8 \\
\bottomrule
\end{tabular}
\end{center}
}

To obtain these counts, split the even-weight \(x\)-masks in \(p^\perp\) by weight. The \(x=0\) quotient contributes \(1,4,3\) classes of costs \(0,1,2\); each of the six weight-two \(x\)-masks contributes four classes of cost \(2\) and four of cost \(3\); the weight-four mask contributes eight classes of cost \(4\).

For every suffix representative \(w\in\xi\), a prefix representative \(v\) of \(e(\xi)\) makes \(v+w\) commute with \(S\). Conversely, commutation forces this prefix class for a given suffix class. Thus the mixed centralizer fiber indexed by \(\xi\) has minimum weight \(d(\xi)+\operatorname{cost}(e(\xi))\). If \(\xi\ne0\), its suffix is nonzero, so no vector in this fiber is \(p\). Each vector in the fiber therefore has weight at least five: if it lies in \(S\), Sections~4--5 leave \(p\) as the only nonzero stabilizer word of weight at most four; otherwise its weight is at least \(d(S)\). Hence
\[
d(\xi)+\operatorname{cost}(e(\xi))\geq5
\qquad(\xi\ne0).
\]

Since \(\xi\mapsto e(\xi)\) is a bijection, the cost distribution gives the following bounds whatever the map \(\tau\) is:

\[
\begin{aligned}
\#\{\xi \ne 0 : d(\xi) &\leq 1\} \leq 8, \\
\#\{\xi \ne 0 : d(\xi) &\leq 2\} \leq 32, \\
\#\{\xi \ne 0 : d(\xi) &\leq 3\} \leq 59.
\end{aligned}
\]

The right-hand sides are the numbers of classes of cost at least 4, 3, and 2, respectively. A code \(C\) that violates any of the three bounds cannot arise from Lemma~\ref{lem:ten-site}.

\section{Exhaustive additive-code classification}\label{exhaustive-additive-code-classification}

\subsection{What is classified}\label{what-is-classified}

Under the standard correspondence between Pauli operators and additive codes \cite{calderbank-shor-sloane}, a quaternary symbol is a pair of bits. A length-\(n\) additive code is a binary subspace of \(\mathbb{F}_{2}^{2n}\), with quaternary Hamming weight equal to physical Pauli weight. Equivalence is coordinate permutation plus an arbitrary local \(\operatorname{GL}(2,2)\) map on each two-bit symbol plane. No \(\operatorname{GF}(4)\)-linearity is assumed. Since \(\operatorname{GL}(2,2)=\operatorname{Sp}(2,2)\), every such map preserves the symplectic form and Pauli weight. Hence the hull dimension of \(C\) and the minimum-weight spectrum of the cosets \(J^\perp_W/C\) used in Section~6.2 are invariants of the equivalence class, so it suffices to check one representative per class.

Grassl, Krotov, Sok, and Solé \cite{grassl-et-al} report 37 equivalence classes in the \((n,k)=(10,10)\) cell of their additive \((n,2^{k}, \geq 5)_{4}\) classification (Section~2.3, Table 1 of \emph{The punctured dodecacode is unique}). They report the count but not the generator matrices, so we reconstruct the classes by lengthening. The census does not use the published counts; they serve only as a cross-check.

\subsection{Why the lengthening census is exhaustive for this cell}\label{why-the-lengthening-census-is-exhaustive-for-this-cell}

The census starts from the zero code in the \((n,k)=(4,0)\) cell and adds one coordinate at a time. The following argument shows that every code in the target cell is reached.

Let \(D\leq\mathbb{F}_2^{2(m+1)}\) be a target code of binary dimension \(k\), and project it onto its last two-bit coordinate. The image has binary dimension \(r\in\{0,1,2\}\); the kernel is a length-\(m\) parent code of binary dimension \(k-r\). Choose a basis of the image and lift its basis vectors to \(D\). Their old-coordinate parts determine \(r\) quotient classes modulo the parent. For \(r=1\), the nonzero class must have minimum old weight at least four. For \(r=2\), each of the three nonzero combinations of the two classes must have minimum old weight at least four. After a local \(\operatorname{GL}(2,2)\) change of basis, these are exactly the extensions below. Conversely, each listed extension has distance at least five whenever its parent has distance at least five and the stated coset minima are at least four.

Thus lengthening a parent of length \(m\) and binary dimension \(k-r\) uses one new coordinate of rank \(r\):

\begin{itemize}
\tightlist
\item
  \(r=0\): append a zero coordinate;
\item
  \(r=1\): choose one quotient coset of the parent whose old minimum weight is at least four and attach a nonzero new symbol;
\item
  \(r=2\): choose two quotient labels whose three nonzero XOR combinations all have old minimum weight at least four and attach independent new symbols.
\end{itemize}

The threshold is four: a word with nonzero new symbol gains one physical unit of weight, while a word with zero new symbol lies in the distance-five parent. Thus every target code appears in one of these three cases.

\Needspace{12\baselineskip}
Table~\ref{tab:predecessors} lists the predecessor cells needed to reach \((10,10)\) and the reconstructed class counts:

{
\begin{longtable}[]{@{}rl@{}}
\caption{Predecessor cells required by the lengthening recurrence to reach \((10,10)\), with reconstructed class counts.}
\label{tab:predecessors}\\
\toprule\noalign{}
length \(n\) & required binary dimensions \(k\): number of classes \\
\midrule\noalign{}
\endhead
\bottomrule\noalign{}
\endlastfoot
5 & \(0:1, 1:1, 2:1, 3:0, 4:0\) \\
6 & \(2:5, 3:1, 4:0, 5:0, 6:0\) \\
7 & \(4:43, 5:1, 6:0, 7:0, 8:0\) \\
8 & \(6:1579, 7:0, 8:0, 9:0, 10:0\) \\
9 & \(8:2298, 9:0, 10:0\) \\
10 & \(10:37\) \\
\end{longtable}
}

The cells with zero classes are computed and found to be empty, not skipped. The recurrence also refers to cells of larger dimension that are not listed, such as \((7,9)\) and \((6,8)\). These are empty as well: a code of binary dimension \(k\) and distance at least five contains a subcode of dimension \(k-1\) with the same property, so emptiness of \((n,k)\) implies emptiness of \((n,k')\) for every \(k'>k\), and each unlisted cell lies above a listed empty cell of the same length. Hence the census reaches every code in the \((10,10)\) cell.

\subsection{Equivalence certificates}\label{equivalence-certificates}

For every candidate the census constructs a colored incidence graph with one vertex for every codeword, three nonzero-symbol vertices for every coordinate, and one coordinate-group vertex binding each triple of symbols. The colors separate the three vertex types. A color-preserving graph isomorphism is exactly a reordering of codewords, a coordinate permutation, and an arbitrary permutation of the three nonzero symbols in each coordinate. The latter is \(\operatorname{GL}(2,2)\). Thus the graph certificate gives a canonical additive-monomial equivalence test; the weight and pair-rank invariants only filter candidates before canonicalization. The graph-isomorphism implementation is based on nauty \cite{mckay-piperno}, via \texttt{pynauty} 2.8.8.1.

The census keeps one representative for every certificate it generates, so no class reached by an extension is discarded. To check that no class was lost or counted twice, we recompute the census with a different graph encoding (\texttt{verification/s7\_mass\_check.py}) and verify the orbit--stabilizer identity in each of the 36 cells \((n,k)\) with \(5\leq n\leq10\) and \(\max(0,2n-10)\leq k\leq10\), including the unlisted empty cells:
\[
\sum_{P}\frac{m!\,6^{m}}{\lvert\operatorname{Aut}P\rvert}\,\operatorname{ext}_r(P)
=\sum_{D}\frac{(m+1)!\,6^{m+1}}{\lvert\operatorname{Aut}D\rvert},
\]
where \(P\) runs over parent classes, \(\operatorname{ext}_r(P)\) counts labelled extensions of rank \(r\), and \(D\) runs over the classes found. Both sides count the labelled codes in the cell, so equality shows that no class is missing or repeated; this is the standard consistency check for classifications \cite{kaski-ostergard}. In the target cell both sides equal \(2{,}333{,}060{,}795{,}842{,}560\). The recomputation finds the same 37 canonical certificates. Table~\ref{tab:aut} lists the automorphism group orders.

{\small
\setlength{\tabcolsep}{5pt}
\begin{center}
\captionof{table}{Hull dimension and automorphism-group order \(\lvert\operatorname{Aut}\rvert\) (coordinate permutations and per-site \(\operatorname{GL}(2,2)\)) of the 37 length-ten classes, numbered as in the repository.}
\label{tab:aut}
\begin{tabular}{@{}rrr@{\hspace{14pt}}rrr@{\hspace{14pt}}rrr@{\hspace{14pt}}rrr@{}}
\toprule
class & hull & \(\lvert\operatorname{Aut}\rvert\) & class & hull & \(\lvert\operatorname{Aut}\rvert\) & class & hull & \(\lvert\operatorname{Aut}\rvert\) & class & hull & \(\lvert\operatorname{Aut}\rvert\) \\
\midrule
1 & 2 & 2 & 2 & 2 & 1 & 3 & 0 & 2 & 4 & 0 & 10 \\
5 & 2 & 2 & 6 & 0 & 2 & 7 & 0 & 2 & 8 & 0 & 30 \\
9 & 4 & 48 & 10 & 2 & 30 & 11 & 2 & 24 & 12 & 4 & 16 \\
13 & 2 & 2 & 14 & 4 & 12 & 15 & 4 & 6 & 16 & 0 & 2 \\
17 & 4 & 16 & 18 & 2 & 8 & 19 & 2 & 1 & 20 & 4 & 5 \\
21 & 4 & 16 & 22 & 4 & 1 & 23 & 2 & 2 & 24 & 8 & 6 \\
25 & 0 & 2 & 26 & 4 & 5 & 27 & 2 & 2 & 28 & 2 & 5 \\
29 & 6 & 5 & 30 & 6 & 10 & 31 & 2 & 5 & 32 & 8 & 27 \\
33 & 2 & 8 & 34 & 4 & 8 & 35 & 4 & 16 & 36 & 2 & 10 \\
37 & 4 & 8 &  &  &  &  &  &  &  &  &  \\
\bottomrule
\end{tabular}
\end{center}
}

\subsection{Hull and coset-capacity results}\label{hull-and-coset-capacity-results}

The census produces 37 generator matrices, each generating a code of size 1,024 and minimum distance five. Separate scripts recompute their ranks, minimum distances, hull dimensions, and coset minima. Their hull dimensions are as follows.

\Needspace{12\baselineskip}
{
\begin{longtable}[]{@{}rr@{}}
\caption{Binary symplectic hull dimensions of the 37 length-ten classes.}
\label{tab:hull}\\
\toprule\noalign{}
binary symplectic hull dimension & number of classes \\
\midrule\noalign{}
\endhead
\bottomrule\noalign{}
\endlastfoot
0 & 7 \\
2 & 14 \\
4 & 12 \\
6 & 2 \\
8 & 2 \\
\end{longtable}
}

\Needspace{17\baselineskip}
By Table~\ref{tab:hull}, the 25 classes with hull dimension different from four cannot arise from the code \(C\) required by Lemma~\ref{lem:ten-site}. For the 12 hull-four classes, the coset-capacity audit gives Table~\ref{tab:coset-capacity}:

{\small
\setlength{\tabcolsep}{4pt}
\captionof{table}{Coset-capacity audit of the 12 hull-four classes. Every class exceeds the capacity \(N_{\leq3}\leq59\) of Section~6.2.}
\label{tab:coset-capacity}
\begin{tabularx}{\linewidth}{@{}rXrrrl@{}}
\toprule
class & suffix weight spectrum \(\{w:N_w\}\) & \(N_{\leq1}\) & \(N_{\leq2}\) & \(N_{\leq3}\) & capacity \\
\midrule
9 & \(\{0:1,2:39,3:24\}\) & 0 & 39 & 63 & fails \\
12 & \(\{0:1,1:2,2:37,3:24\}\) & 2 & 39 & 63 & fails \\
14 & \(\{0:1,2:39,3:24\}\) & 0 & 39 & 63 & fails \\
15 & \(\{0:1,2:27,3:36\}\) & 0 & 27 & 63 & fails \\
17 & \(\{0:1,1:2,2:37,3:24\}\) & 2 & 39 & 63 & fails \\
20 & \(\{0:1,2:35,3:27,4:1\}\) & 0 & 35 & 62 & fails \\
21 & \(\{0:1,1:2,2:37,3:24\}\) & 2 & 39 & 63 & fails \\
22 & \(\{0:1,2:23,3:40\}\) & 0 & 23 & 63 & fails \\
26 & \(\{0:1,2:35,3:27,4:1\}\) & 0 & 35 & 62 & fails \\
34 & \(\{0:1,1:2,2:37,3:24\}\) & 2 & 39 & 63 & fails \\
35 & \(\{0:1,1:2,2:37,3:24\}\) & 2 & 39 & 63 & fails \\
37 & \(\{0:1,1:2,2:37,3:24\}\) & 2 & 39 & 63 & fails \\
\bottomrule
\end{tabularx}
}

In every row \(N_{\leq3}\in\{62,63\}>59\). Eight rows also have \(N_{\leq2}=39>32\).

\begin{lemma}[Additive lift exclusion]\label{lem:additive-lift}
None of the 37 length-ten classes can supply the code in Lemma~\ref{lem:ten-site}: 25 have hull dimension different from four, and the remaining 12 violate the coset-capacity bound.
\end{lemma}
\begin{proof}
The first exclusion follows from Table~\ref{tab:hull}. For every remaining class, Table~\ref{tab:coset-capacity} gives \(N_{\leq3}\in\{62,63\}\), exceeding the necessary capacity 59. Hence none of the 37 classes can supply the required code.
\end{proof}

\section{Proof of the main theorem}\label{discussion-and-conclusion}

\begin{proof}[Proof of Theorem~\ref{thm:main}]
Assume a binary \(\left[\!\left[14,3,d\geq 5\right]\!\right]\) stabilizer code exists. Sections~3--4 give \(A_{1}=A_{3}=0\), \(A_{2}\leq2\), and \(A_{4}\in\{1,3\}\). Section~5.3 then excludes weight-two stabilizer words, so \(A_{2}=0\). Sections~5.3--5.5 exclude all pairs of distinct weight-four stabilizer words whose supports intersect in at least two sites and classify the \(A_4=3\) case into five support geometries. Section~5.6 excludes the four overlapping rank-three geometries, and Section~5.7 excludes the disjoint \(4+4+4+2\) geometry. Hence \(A_{4}=1\), which is Lemma~\ref{lem:unique-check}.

Lemma~\ref{lem:ten-site} then forces a ten-site additive code of size 1,024, distance at least five, and hull dimension four. By Lemma~\ref{lem:additive-lift}, none of the 37 equivalence classes of length-ten additive codes of size 1,024 and distance at least five can supply the code required by Lemma~\ref{lem:ten-site}: 25 have the wrong hull dimension, and the remaining 12 violate the coset-capacity bound. This is a contradiction.
\end{proof}

The argument is specific to these parameters and to stabilizer codes; it gives no information about other lengths or about non-stabilizer codes.

\section{Computational material}\label{computational-material}

The computations are the exact certificates of Sections~4--5 and the classification of Section~7. All arithmetic is exact, and no floating-point computation is used. The code and data are at
\begin{center}
\url{https://github.com/brandonlign/quantum-code-bounds}
\end{center}
With Python 3.10 or newer and Node.js, the command
\begin{center}
\texttt{bash verification/replay.sh}
\end{center}
checks every certificate of Sections~4--7, most of them in both Python and JavaScript (with BigInt arithmetic), and checks the 37 stored class representatives. The option \texttt{--census} also recomputes the classification and runs the mass check of Section~\ref{equivalence-certificates}; this takes a few minutes and requires \texttt{pynauty==2.8.8.1}. Script names follow the section numbers of this paper.

Every Farkas certificate of Section~5 is also provided as a data file in \texttt{verification/certificates/}, in the format of Appendix~\ref{app:systems}. A short standalone program, \texttt{verification/check\_certificates.py}, rebuilds each system from the data file and the definitions of Appendix~\ref{app:systems} alone. A reader can therefore check Section~5 without the rest of the code, or write an independent checker.

\section*{Acknowledgments}

Generative AI tools (ChatGPT, Codex and Claude) were used to explore proof strategies, check derivations, write verification code, and edit the text. The author directed the research, checked the arguments and computations, and takes full responsibility for the content.

\section*{Statements and declarations}

\textbf{Data availability.} All code, certificate data and the 37 class representatives are publicly available in the repository given in Section~\ref{computational-material}. They are also archived on Zenodo at \href{https://doi.org/10.5281/zenodo.22885774}{doi:10.5281/zenodo.22885774}.

\textbf{Funding.} No funding was received for this work.

\textbf{Competing interests.} The author has no competing interests to declare.

\Needspace{20\baselineskip}
\begin{thebibliography}{99}
\bibitem{calderbank-shor-sloane}
A.~R. Calderbank, E.~M. Rains, P.~W. Shor, and N.~J.~A. Sloane,
``Quantum error correction via codes over GF(4),''
\emph{IEEE Trans. Inf. Theory} 44 (1998), 1369--1387.
\href{https://doi.org/10.1109/18.681315}{doi:10.1109/18.681315}.

\bibitem{rains-weight}
E.~M. Rains, ``Quantum weight enumerators,''
\emph{IEEE Trans. Inf. Theory} 44 (1998), 1388--1394.
\href{https://doi.org/10.1109/18.681316}{doi:10.1109/18.681316}.

\bibitem{rains-shadow}
E.~M. Rains, ``Quantum shadow enumerators,''
\emph{IEEE Trans. Inf. Theory} 45 (1999), 2361--2366.
\href{https://doi.org/10.1109/18.796376}{doi:10.1109/18.796376}.

\bibitem{ball-centelles-huber}
S. Ball, A. Centelles, and F. Huber, ``Quantum error-correcting codes and their geometries,''
\emph{Ann. Inst. H. Poincaré D} 10 (2023), no.~2, 337--405.
\href{https://doi.org/10.4171/AIHPD/160}{doi:10.4171/AIHPD/160}.

\bibitem{bierbrauer-et-al}
J. Bierbrauer, R. Fears, S. Marcugini, and F. Pambianco,
``The nonexistence of a [[13,5,4]]-quantum stabilizer code,''
\emph{IEEE Trans. Inf. Theory} 57 (2011), 4788--4793.
\href{https://doi.org/10.1109/TIT.2011.2146430}{doi:10.1109/TIT.2011.2146430}.

\bibitem{grassl-et-al}
M. Grassl, D. Krotov, L. Sok, and P. Solé,
``The punctured dodecacode is unique,''
\emph{Mathematics and Education in Mathematics} 55 (2026), 393--404.
\href{https://doi.org/10.55630/mem.2026.55.393-404}{doi:10.55630/mem.2026.55.393-404}.

\bibitem{kaski-ostergard}
P. Kaski and P.~R.~J. \"Osterg\aa rd,
\emph{Classification Algorithms for Codes and Designs},
Springer, Berlin, 2006.
\href{https://doi.org/10.1007/3-540-28991-7}{doi:10.1007/3-540-28991-7}.

\bibitem{mckay-piperno}
B.~D. McKay and A. Piperno, ``Practical graph isomorphism, II,''
\emph{J. Symbolic Comput.} 60 (2014), 94--112.
\href{https://doi.org/10.1016/j.jsc.2013.09.003}{doi:10.1016/j.jsc.2013.09.003}.

\bibitem{grassl-table}
M. Grassl, ``Bounds on the minimum distance of additive quantum codes, (n=14,k=3,q=4),''
maintained code table, accessed 24 September 2026.
\href{https://codetables.de/QECC.php?k=3&n=14&q=4}{codetables.de}.
\end{thebibliography}

\appendix
\section{Exact specification of the Section 5 systems}\label{app:systems}

This appendix defines every linear system of Section~5 in enough detail that a reader can rebuild it, and check its certificate, without reference to any code supplied by the author.

\subsection{Variables}\label{app:variables}

The fourteen qubits are split into consecutive blocks of sizes \(m_1,\dots,m_r\) and a suffix of size \(f\), as listed in Table~\ref{tab:app-systems}. For a block \(i\) with subgroup \(R_i\leq V_{m_i}\), compute the normalizer \(R_i^\perp\cap V_{m_i}\) and split it into complete \(R_i\)-cosets. Each coset \(Q\) has a weight histogram \(p(Q)=(p_0,\dots,p_{m_i})\) as in Section~5.1. The distinct histograms, sorted lexicographically, are the \emph{patterns} \(P_{i,1},P_{i,2},\dots\) of block \(i\).

A column is a tuple \(y=(k_1,\dots,k_r,b)\): one pattern index for each block and a suffix weight \(0\leq b\leq f\). Columns are ordered lexicographically. Its variable counts the \(R\)-cosets of \(S\) (where \(R=R_1\oplus\dots\oplus R_r\)) whose block histograms are \(P_{1,k_1},\dots,P_{r,k_r}\) and whose suffix part has weight \(b\). In the two systems marked ``cells'' in Table~\ref{tab:app-systems}, the variables are instead the split counts \(H_{a,b}\) themselves, which is the same as taking the patterns of block 1 to be the unit vectors \(e_0,\dots,e_{m_1}\).

A split cell is \(c=(a_1,\dots,a_r,a_f)\) with \(0\leq a_i\leq m_i\), \(0\leq a_f\leq f\), and \(|c|=\sum a_i+a_f\). The contribution of column \(y\) to the three enumerators in cell \(c\) is
\[
\begin{aligned}
H_c(y)&=\prod_{i=1}^r P_{i,k_i}[a_i]\cdot[a_f=b],\\
T_c(y)&=\prod_{i=1}^r\Bigl(\sum_{j}K^{m_i}_{a_i}(j)\,P_{i,k_i}[j]\Bigr)\cdot K^{f}_{a_f}(b),\\
\mathtt{SH}_c(y)&=\prod_{i=1}^r\Bigl(\sum_{j}(-1)^jK^{m_i}_{a_i}(j)\,P_{i,k_i}[j]\Bigr)\cdot(-1)^bK^{f}_{a_f}(b),
\end{aligned}
\]
which is the tensor-product form of the transforms in Section~3. Here \(T_c=2048\,C_c\) and \(\mathtt{SH}_c=2048\,\Sh_c\).

\subsection{Rows}\label{app:rows}

Each constraint is identified by a type and, where relevant, a cell. Table~\ref{tab:app-rows} gives the coefficient of column \(y\) in each type. Every row holds for every code containing \(R\), by Sections~3--5.1.

\Needspace{16\baselineskip}
{\small
\begin{center}
\captionof{table}{Row types. Equalities (\(=\)) take free-sign multipliers \(\nu\); inequalities (\(\leq\)) take multipliers \(\lambda\geq0\).}
\label{tab:app-rows}
\begin{tabularx}{\linewidth}{@{}lXl@{}}
\toprule
type & coefficient of column \(y\) & relation \\
\midrule
\texttt{unit} & \(H_{(0,\dots,0)}(y)\) & \(=1\) \\
\texttt{total} & \(\sum_c H_c(y)\) & \(=2048\) \\
\texttt{low} \(c\), \(1\leq|c|\leq4\) & \(T_c(y)-2048\,H_c(y)\) & \(=0\) \\
\texttt{zero} \(c\), \(|c|\in\{1,3\}\) & \(H_c(y)\) & \(=0\) \\
\texttt{weight} \(w\) & \(\sum_{|c|=w}H_c(y)\) & \(=A_w\) \\
\texttt{coset\_eq}, \texttt{coset\_le} & \(\sum_a \omega_a H_{(a,b)}(y)\) for given \(\omega,b\) & \(=0\), \(\leq0\) \\
\texttt{incl} \(c\) & \(2048\,H_c(y)-T_c(y)\) & \(\leq0\) \\
\texttt{cnonneg} \(c\) & \(-T_c(y)\) & \(\leq0\) \\
\texttt{shadow} \(c\) & \(-\mathtt{SH}_c(y)\) & \(\leq0\) \\
\bottomrule
\end{tabularx}
\end{center}
}

The \texttt{weight} rows use only the values \(A_1=A_2=A_3=0\) and the case assumption \(A_4=3\), and appear only in Sections~5.6--5.7. The \texttt{zero} row uses \(A_1=0\) from Lemma~\ref{lem:low-weight} and appears only in Section~5.4.2. The \texttt{coset} rows appear only in the two ``cells'' systems, where they encode the coset relations stated in Sections~5.3 and~5.4.2: \(H_{2,b}\geq H_{0,b}\), \(H_{1,b}=0\), and \(H_{4,b}=H_{2,b}+3H_{0,b}\). Each such relation is implied by every coset pattern of its subgroup, which a checker should confirm.

\subsection{Certificates}\label{app:certificates}

A certificate is a finite list of (row, multiplier) pairs. For each system, the residual \(c(y)=\sum\lambda_\rho\,G_\rho(y)+\sum\nu_\rho\,M_\rho(y)\) must satisfy \(c(y)\geq0\) in every column, with \(\nu\cdot\rhs<0\). Section~5.2 shows that no code satisfies the system in that case. Table~\ref{tab:app-systems} lists all eleven systems. The multipliers themselves are in \texttt{verification/certificates/}, one JSON file per row of the table. Each file contains the block sizes, the subgroup generators, the coset patterns with multiplicities, and the labelled rows with their integer multipliers. The certificates are sparse: they have between 22 and 111 nonzero multipliers.

\Needspace{18\baselineskip}
{\small
\setlength{\tabcolsep}{3pt}
\begin{center}
\captionof{table}{The eleven Section~5 systems. Blocks give \((m_1,\dots,m_r;f)\). ``Pat.'' is the number of coset patterns per block, and ``rows used'' is the number of rows, equalities and inequalities, with a nonzero multiplier. Generators act on the block sites in order, and in Section~5.7 each of the three blocks carries \(\mathtt{ZZZZ}\).}
\label{tab:app-systems}
\begin{tabularx}{\linewidth}{@{}l>{\raggedright\arraybackslash}Xlrrrr@{}}
\toprule
\S & generators of \(R\) & blocks & pat. & columns & rows used & \(\nu\cdot\rhs\) \\
\midrule
5.3 & \(\mathtt{ZZ}\) & \((2;12)\) & cells & 39 & 23 & \(-2571108352\) \\
5.4.1 & \(\mathtt{ZZZ}\) & \((3;11)\) & 3 & 36 & 22 & \(-22528\) \\
5.4.2 & \(\mathtt{XXXX},\mathtt{ZZZZ}\) & \((4;10)\) & cells & 55 & 39 & \(-320017816092672\) \\
5.4.3 & \(\mathtt{XXZZI},\mathtt{ZZZIZ}\) & \((5;9)\) & 9 & 90 & 38 & \(-52501014528\) \\
5.4.4 & \(\mathtt{ZZZZII},\mathtt{ZZIIZZ}\) & \((6;8)\) & 13 & 117 & 25 & \(-8820736\) \\
5.4.4 & \(\mathtt{ZZZZII},\mathtt{XXIIXX}\) & \((6;8)\) & 14 & 126 & 27 & \(-22212608\) \\
5.6 & three distinct intersections & \((9;5)\) & 40 & 240 & 32 & \(-86470656\) \\
5.6 & one common site & \((10;4)\) & 27 & 135 & 33 & \(-9031680\) \\
5.6 & two intersections & \((10;4)\) & 60 & 300 & 34 & \(-20484096\) \\
5.6 & one intersection & \((11;3)\) & 74 & 296 & 33 & \(-61440\) \\
5.7 & \(\mathtt{ZZZZ}\) on each block & \((4,4,4;2)\) & 6,6,6 & 648 & 111 & \(-125829120\) \\
\bottomrule
\end{tabularx}
\end{center}
}

The Section~5.6 generators are the \(\mathtt{Z}\)-only words on the supports in Table~\ref{tab:support-geometry}. The patterns of the 5.4.1, 5.4.2 and 5.7 blocks appear in Sections~5.4.1, 5.4.2 and Table~\ref{tab:block-patterns}. For the remaining two-generator blocks, the patterns (written \((p_0,\dots,p_m)\!:\)multiplicity) are:

{\raggedright
\begin{itemize}
\item \(\langle\mathtt{XXZZI},\mathtt{ZZZIZ}\rangle\): \mbox{\((0,0,0,0,4,0)\!:\!12\)}, \mbox{\((0,0,0,1,0,3)\!:\!8\)}, \mbox{\((0,0,0,2,0,2)\!:\!12\)}, \mbox{\((0,0,0,3,0,1)\!:\!12\)}, \mbox{\((0,0,0,4,0,0)\!:\!1\)}, \mbox{\((0,0,1,0,3,0)\!:\!12\)}, \mbox{\((0,0,2,0,2,0)\!:\!3\)}, \mbox{\((0,1,0,2,0,1)\!:\!3\)}, \mbox{\((1,0,0,0,3,0)\!:\!1\)}.
\item \(\langle\mathtt{ZZZZII},\mathtt{ZZIIZZ}\rangle\): \mbox{\((0,0,0,0,0,0,4)\!:\!16\)}, \mbox{\((0,0,0,0,0,4,0)\!:\!24\)}, \mbox{\((0,0,0,0,2,0,2)\!:\!24\)}, \mbox{\((0,0,0,0,3,0,1)\!:\!64\)}, \mbox{\((0,0,0,0,4,0,0)\!:\!12\)}, \mbox{\((0,0,0,1,0,3,0)\!:\!64\)}, \mbox{\((0,0,0,2,0,2,0)\!:\!24\)}, \mbox{\((0,0,0,4,0,0,0)\!:\!2\)}, \mbox{\((0,0,1,0,2,0,1)\!:\!12\)}, \mbox{\((0,0,2,0,2,0,0)\!:\!6\)}, \mbox{\((0,0,3,0,0,0,1)\!:\!1\)}, \mbox{\((0,1,0,2,0,1,0)\!:\!6\)}, \mbox{\((1,0,0,0,3,0,0)\!:\!1\)}.
\item \(\langle\mathtt{ZZZZII},\mathtt{XXIIXX}\rangle\): \mbox{\((0,0,0,0,0,4,0)\!:\!32\)}, \mbox{\((0,0,0,0,1,0,3)\!:\!16\)}, \mbox{\((0,0,0,0,2,0,2)\!:\!32\)}, \mbox{\((0,0,0,0,3,0,1)\!:\!40\)}, \mbox{\((0,0,0,0,4,0,0)\!:\!16\)}, \mbox{\((0,0,0,1,0,3,0)\!:\!48\)}, \mbox{\((0,0,0,2,0,2,0)\!:\!32\)}, \mbox{\((0,0,0,3,0,1,0)\!:\!4\)}, \mbox{\((0,0,1,0,1,0,2)\!:\!8\)}, \mbox{\((0,0,1,0,2,0,1)\!:\!16\)}, \mbox{\((0,0,1,0,3,0,0)\!:\!5\)}, \mbox{\((0,0,2,0,1,0,1)\!:\!2\)}, \mbox{\((0,1,0,1,0,2,0)\!:\!4\)}, \mbox{\((1,0,0,0,2,0,1)\!:\!1\)}.
\end{itemize}\par}

The 40--74 patterns of the Section~5.6 blocks are too long to print, but they are determined by the generators. They are listed in the certificate files, and \texttt{check\_certificates.py} recomputes them before using them.

\end{document}
